% Part: methods
% Chapter: induction
% Section: structural-induction

\documentclass[../../../include/open-logic-section]{subfiles}

\begin{document}

\olfileid{mth}{ind}{sti}

\olsection{સંરચનાત્મક આગમન}

અત્યાર સુધી આપણે બધી પ્રાકૃતિક સંખ્યાઓ વિશેનાં
પરિણામો સ્થાપવા આગમન વાપર્યું છે. પરંતુ અનુરૂપ
સિદ્ધાંત આગમનાત્મક રીતે વ્યાખ્યાયિત ગણના બધા
!!{element}s વિશે પરિણામો સાબિત કરવા સીધો વાપરી
શકાય છે. તેને ઘણી વાર \emph{સંરચનાત્મક} આગમન
કહે છે, કારણ કે તે આગમનાત્મક રીતે વ્યાખ્યાયિત
પદાર્થોની સંરચના પર આધાર રાખે છે.

સામાન્ય રીતે આગમનાત્મક વ્યાખ્યામાં (a)~ગણના
"શરૂઆતના" !!{element}sની યાદી અને (b)~ગણના
અગાઉના ઘટકોમાંથી નવા !!{element}s
બનાવતી ક્રિયાઓની યાદી હોય છે. દાખલા તરીકે,
સુઘડ પદોના કિસ્સામાં શરૂઆતના પદાર્થો અક્ષરો છે.
એક જ ક્રિયા છે:
\begin{align*}
  o(s_1, s_2) = & [s_1 \circ s_2]
\end{align*}
પ્રાકૃતિક સંખ્યાઓ~$\Nat$ને પણ આગમનાત્મક
વ્યાખ્યાથી મળેલી માની શકો: શરૂઆતનો પદાર્થ~$0$
છે અને ક્રિયા અનુગામી વિધેય~$x + 1$ છે.

આગમનાત્મક રીતે વ્યાખ્યાયિત ગણના બધા ઘટકો
વિશે કંઈક સાબિત કરવા, એટલે ગણના દરેક
!!{element}માં ગુણધર્મ~$P$ છે એમ બતાવવા,
આપણે:
\begin{enumerate}
\item શરૂઆતના પદાર્થોમાં~$P$ છે એમ સાબિત કરવું, અને
\item દરેક ક્રિયા~$o$ માટે, તેના નિવેશોમાં~$P$
  હોય તો પરિણામ માટે પણ એવું જ છે એમ સાબિત કરવું.
\end{enumerate}
દાખલા તરીકે, બધાં સુઘડ પદો વિશે કંઈક સાબિત
કરવા પહેલાં બધા અક્ષરો માટે તે સાચું છે એમ
અને પછી $s_1$ અને $s_2$ બંને માટે તે સાચું
હોય તો $[s_1 \circ s_2]$ માટે પણ સાચું છે
એમ સાબિત કરીએ.

\begin{prop}
  કોઈ પણ સુઘડ પદ~$t$માં $[$ની સંખ્યા $]$ની
  સંખ્યા જેટલી જ છે.
\end{prop}

\begin{proof}
સંરચનાત્મક આગમન વાપરીએ. સુઘડ પદો આગમનાત્મક
રીતે વ્યાખ્યાયિત છે: અક્ષરો શરૂઆતના પદાર્થો છે
અને $o$ ક્રિયા જૂનાં સુઘડ પદોમાંથી નવાં બનાવે છે.
\begin{enumerate}
\item દરેક અક્ષર માટે દાવો સાચો છે, કારણ કે એક
  અક્ષરમાં $[$ની સંખ્યા~$0$ છે અને $]$ની
  સંખ્યા પણ~$0$ છે.
\item ધારો કે $s_1$માં $[$ની સંખ્યા $]$ની
  સંખ્યા જેટલી છે, અને $s_2$ માટે પણ એવું
  જ છે. $o(s_1, s_2)$માં, એટલે
  $[s_1 \circ s_2]$માં, $[$ની સંખ્યા એ
  $s_1$ અને $s_2$માંના $[$ની સંખ્યાનો
  સરવાળો વત્તા એક છે. $o(s_1, s_2)$માં
  $]$ની સંખ્યા એ $s_1$ અને $s_2$માંના
  $]$ની સંખ્યાનો સરવાળો વત્તા એક છે.
  તેથી $o(s_1, s_2)$માં $[$ની સંખ્યા
  $o(s_1,s_2)$માં $]$ની સંખ્યા જેટલી છે.
\end{enumerate}
\end{proof}

\begin{prob}
  કોઈ પણ સુઘડ પદની શરૂઆત~$]$થી થતી નથી
  એ સંરચનાત્મક આગમનથી સાબિત કરો.
\end{prob}

સંરચનાત્મક આગમનથી બીજી સાબિતી જોઈએ:
ચિહ્ન-શ્રેણી~$t$નો અરિક્ત ઉચિત આરંભિક ખંડ
એવી શ્રેણી~$s$ છે જે ડાબેથી વાંચતાં ચિહ્ને
ચિહ્ને~$t$ સાથે મળે છે, પરંતુ $t$ વધુ લાંબી
છે. દાખલા તરીકે, $[a \circ {}$ એ
$[a \circ b]$નો ઉચિત આરંભિક ખંડ છે; પરંતુ
$[b \circ {}$ નથી (બીજા ચિહ્ને ભેદ છે)
અને $[a \circ b]$ પણ નથી (લંબાઈ સરખી છે).
\sourcecorrection{OLMD-165}{સ્થિર અંગ્રેજી મૂળમાં ઉચિત આરંભિક
  ખંડ અરિક્ત હોવાની શરત કહેવાઈ નથી. ખાલી શ્રેણી
  પણ ઉચિત આરંભિક ખંડ ગણીએ તો આગળનું પ્રતિજ્ઞાપન
  ખોટું થાય: તેમાં $[$ અને $]$ બંનેની સંખ્યા શૂન્ય છે.
  આગળની સાબિતીમાં ખાલી ખંડનો કિસ્સો નથી; અહીં
  અભિપ્રેત અરિક્ત શરત સ્પષ્ટ કરી છે.}

\begin{prop}\ollabel{prop:initial}
  સુઘડ પદ~$t$ના દરેક અરિક્ત ઉચિત આરંભિક ખંડમાં
  $[$ની સંખ્યા $]$ની સંખ્યાથી વધુ છે.
\end{prop}

\begin{proof}
  $t$ પર આગમન કરીએ:
  \begin{enumerate}
  \item $t$ પોતે એક અક્ષર છે: ત્યારે $t$નો
    કોઈ અરિક્ત ઉચિત આરંભિક ખંડ નથી.
  \item કોઈ સુઘડ પદો $s_1$ અને~$s_2$ માટે
    $t = [s_1 \circ s_2]$ છે. $r$ જો $t$નો
    અરિક્ત ઉચિત આરંભિક ખંડ હોય, તો નીચેની
    શક્યતાઓ છે:
    \begin{enumerate}
    \item $r$ ફક્ત $[$ છે: ત્યારે $r$માં
      $]$ કરતાં એક $[$ વધુ છે.
    \item $r$ એ $[r_1$ છે, જ્યાં $r_1$ એ~$s_1$નો
      અરિક્ત ઉચિત આરંભિક ખંડ છે: $s_1$ સુઘડ પદ
      હોવાથી આગમન-પરિકલ્પના મુજબ $r_1$માં
      $]$ કરતાં વધુ $[$ છે અને $[r_1$માં
      પણ એવું જ છે.
    \item $r$ એ $[s_1$ અથવા $[s_1 \circ {}$
      છે: અગાઉના પરિણામ મુજબ~$s_1$માં $[$
      અને $]$ની સંખ્યા સરખી છે; તેથી
      $[s_1$ અથવા $[s_1 \circ {}$માં
      $]$ કરતાં એક $[$ વધુ છે.
    \item $r$ એ $[s_1 \circ r_2$ છે, જ્યાં
      $r_2$ એ~$s_2$નો અરિક્ત ઉચિત આરંભિક ખંડ છે:
      આગમન-પરિકલ્પના મુજબ $r_2$માં $]$
      કરતાં વધુ $[$ છે. અગાઉના પરિણામ
      મુજબ~$s_1$માં $[$ અને $]$ની સંખ્યા
      સરખી છે. તેથી $[s_1 \circ r_2$માં
      $]$ કરતાં વધુ $[$ છે.
    \item $r$ એ $[s_1 \circ s_2$ છે:
      અગાઉના પરિણામ મુજબ $s_1$માં $[$
      અને $]$ની સંખ્યા સરખી છે, અને~$s_2$
      માટે પણ એવું જ છે. તેથી
      $[s_1 \circ s_2$માં $]$ કરતાં એક
      $[$ વધુ છે.
    \end{enumerate}
  \end{enumerate}
\end{proof}

\end{document}
